\documentclass[11pt]{article}
\usepackage[T1]{fontenc}
\usepackage[utf8]{inputenc}
\usepackage{lmodern}
\usepackage{microtype}
\usepackage[margin=1in]{geometry}
\usepackage{amsmath,amssymb,amsthm}
\usepackage{booktabs}
\usepackage{array}
\usepackage{longtable}
\usepackage{hyperref}
\usepackage{xcolor}

\hypersetup{
  colorlinks=true,
  linkcolor=blue,
  urlcolor=blue,
  citecolor=blue
}

\setlength{\parskip}{0.6em}
\setlength{\parindent}{0pt}
\setlength{\LTpre}{0.5em}
\setlength{\LTpost}{0.5em}

\newcommand{\pp}[1]{\left[#1\right]_+}
\newcommand{\ii}{\mathrm{i}}

\begin{document}

\begin{center}
{\LARGE Rewritten Thinking Log for the Two-Minus Water-Wave Benchmark}
\end{center}

\begin{abstract}
This document contains a detailed rewritten summary of the reasoning process
used to analyze the two-minus sector benchmark in
\texttt{waterhedron\_benchmark\_blind/case\_1}, extract a closed formula from
Berends--Giele data, and verify it against the supplied implementation. The
goal here is completeness and chronology: I record the main hypotheses, the
failed directions, the sample data that changed my mind, the intermediate local
formulas, and the final exact verification checks. This is a cleaned-up process
narrative rather than a verbatim hidden chain-of-thought dump.
\end{abstract}

\section*{Rewritten Thinking Log}

\subsection*{1. Starting point and initial expectation}

The prompt made a strong claim: in the two-minus sector
\[
\sigma = (-1,-1,+1,\ldots,+1),
\]
the amplitude should be a single analytic formula valid for all $n \ge 4$, and
the hint explicitly described the target as a rational function assembled from
physical factorization-channel poles. So my initial working assumption was:

\medskip
\textbf{Hypothesis A.} The supplied Berends--Giele recursion should reveal a
single global rational expression after enough low-point data are collected and
fitted.
\medskip

The fastest way to test that assumption was to start at the smallest relevant
multiplicity and work upward only if necessary.

\subsection*{2. Isolating the recursion core}

The provided \texttt{OnShellBG.m} contains three conceptually separate parts:

\begin{enumerate}
\item the exact interaction kernels,
\item the exact Berends--Giele recursion and kinematic solver,
\item a long built-in test block.
\end{enumerate}

For fitting and repeated evaluation, the test block was overhead. So the first
practical step was to copy only the recursion and the solver into a minimal
helper file, \texttt{bg\_core.wl}, inside the output directory. That helper
preserved the exact same amplitude definition while letting me evaluate custom
kinematics on demand.

This was not mathematically interesting, but it mattered operationally because
the later search involved many small exact evaluations and repeated restarts.

\subsection*{3. First direct attack: symbolic $n=4$}

The most natural first attempt was a direct symbolic four-point computation.
With the two-minus sign pattern, \texttt{MakeKinematics[4,\{a,b\},sig,1]}
produces four on-shell frequencies and momenta that are rational functions of
$a$ and $b$, so in principle one could ask Mathematica for
\texttt{Together[BGAmplitude[ks,ws,1]]} and factor the result.

That attempt did not give a clean answer. The raw symbolic expression expanded
into a huge formula involving many nested absolute values. More importantly,
concrete rational four-point evaluations repeatedly produced exact
\texttt{Indeterminate} rather than a finite amplitude:
\[
\{2,3\} \mapsto \{-3,2,3,-2\}, \qquad A_4 = \texttt{Indeterminate},
\]
\[
\{3,2\} \mapsto \{-2,3,2,-3\}, \qquad A_4 = \texttt{Indeterminate},
\]
\[
\{5/2,7/2\} \mapsto \{-7/2,5/2,7/2,-5/2\}, \qquad A_4 = \texttt{Indeterminate}.
\]

So the first real obstacle was not fitting; it was that the exact BG recursion
as supplied sits directly on singular four-point internal channels in this
sector. At that stage I considered a stronger tentative possibility:

\medskip
\textbf{Hypothesis B.} Perhaps the prompt's claim of a single global rational
function is simply false already at four points, and the sector is genuinely
piecewise or singular.
\medskip

I did \emph{not} adopt that as the final conclusion, because the right way to
separate ``real amplitude structure'' from ``an artifact of this exact
recursive representation'' was to move to $n \ge 5$, where the code evaluates
cleanly at generic points, and only then circle back to four points.

\subsection*{4. First exact $n=5$ data and homogeneity check}

Once I moved to five points, the recursion became usable immediately. A few
generic exact samples were:
\[
\texttt{freeW}=\{2,5/2,3\}
\quad\Longrightarrow\quad
\omega = \{-9/2,2,5/2,3,-3\},
\quad
A_5=-2304\,\ii,
\]
\[
\texttt{freeW}=\{3/2,2,7/3\}
\quad\Longrightarrow\quad
\omega = \{-53/15,3/2,2,7/3,-23/10\},
\quad
A_5=-\frac{4293}{10}\,\ii,
\]
\[
\texttt{freeW}=\{5/4,7/4,9/4\}
\quad\Longrightarrow\quad
\omega = \{-13/4,5/4,7/4,9/4,-2\},
\quad
A_5=-158.69140625\,\ii.
\]

Two immediate patterns appeared.

First, every amplitude I sampled was purely imaginary. That suggested pulling
out a universal factor of $\ii$ from the start.

Second, scaling all free frequencies by a common factor $c$ gave a clean power
law. Using the seed point $\{2,5/2,3\}$:
\[
c=1 \Longrightarrow A_5=-2304\,\ii,
\]
\[
c=\frac32 \Longrightarrow A_5=-26244\,\ii,
\]
\[
c=2 \Longrightarrow A_5=-147456\,\ii.
\]
Since
\[
\frac{-26244}{-2304}=\left(\frac32\right)^6,
\qquad
\frac{-147456}{-2304}=2^6,
\]
the five-point amplitude is homogeneous of degree $6$ in the frequencies.

That strongly supported:

\medskip
\textbf{Hypothesis C.} After dividing by $\ii$, the two-minus amplitude is a
real homogeneous function of degree $2n-4$.
\medskip

This degree count later matched the final formula exactly.

\subsection*{5. Rational-fit hypothesis on a one-parameter slice}

At this point I tried to follow the prompt's suggested strategy literally. I
chose a one-parameter five-point slice,
\[
\texttt{freeW}=\{x,2,3\},
\]
and sampled exact amplitudes along that line. The resulting values were:

\begin{center}
\small
\begin{tabular}{cccc}
\toprule
$x$ & $\omega_1$ & $\omega_5$ & $A_5$ \\
\midrule
$1$ & $-4$ & $-2$ & $-64\,\ii$ \\
$3/2$ & $-53/13$ & $-63/26$ & $-(12879/26)\,\ii$ \\
$7/4$ & $-37/9$ & $-95/36$ & $-(621859/576)\,\ii$ \\
$9/4$ & $-121/29$ & $-357/116$ & $-(106722/29)\,\ii$ \\
$5/2$ & $-21/5$ & $-33/10$ & $-5712\,\ii$ \\
$11/4$ & $-131/31$ & $-437/124$ & $-(256498/31)\,\ii$ \\
$13/4$ & $-47/11$ & $-175/44$ & $-(10332621/704)\,\ii$ \\
$7/2$ & $-73/17$ & $-143/34$ & $-(584073/34)\,\ii$ \\
$4$ & $-13/3$ & $-14/3$ & $-19968\,\ii$ \\
$9/2$ & $-83/19$ & $-195/38$ & $-(430272/19)\,\ii$ \\
$5$ & $-22/5$ & $-28/5$ & $-25344\,\ii$ \\
\bottomrule
\end{tabular}
\end{center}

The singular values on this slice were also informative:
\[
x=2 \Longrightarrow A_5=\texttt{Indeterminate},
\qquad
x=3 \Longrightarrow A_5=\texttt{Indeterminate}.
\]

Those two singular points looked, at first glance, exactly like the sort of
simple channel poles the prompt had prepared me to expect. So I tested:

\medskip
\textbf{Hypothesis D.} On this slice, $A_5/\ii$ should be a low-degree rational
function of $x$, perhaps with simple poles at $x=2$ and $x=3$.
\medskip

I then tried several increasingly explicit reconstructions:

\begin{enumerate}
\item direct low-degree rational interpolation in SymPy,
\item polynomial interpolation after multiplying by $(x-2)^a(x-3)^b$ for small
integers $a,b$,
\item auxiliary scans with a further factor $(x+5)^p$ to catch the possibility
that the denominator inherited the kinematic solver's $(a+b+c)$-type
structure.
\end{enumerate}

None of those low-degree fits stabilized into a convincing global rational
slice formula. That was the point where the literal rational-fit strategy
stopped looking like the right lens for the data I actually had.

The conclusion was not yet ``the prompt is wrong.'' The more careful
conclusion was:

\medskip
\textbf{Interim conclusion.} The raw slice data do not support an obvious
low-degree rational formula in a single variable, even though singular points
exist on special loci. Any useful closed form is probably hidden by chamber
structure.
\medskip

\subsection*{6. Chamber-fixed symbolic reimplementation}

To get behind that chamber structure, I wrote a second exact implementation,
\texttt{symbolic\_bg.py}. Its purpose was not speed; its purpose was to replace
every absolute value $\lvert\cdot\rvert$ by a fixed sign choice in one generic
ordering chamber and then simplify exactly inside that chamber.

This required a few practical fixes before it became useful:

\begin{enumerate}
\item the local Python environment was older than expected, so the initial
version using \texttt{dataclasses}, modern type annotations, and
\texttt{from \_\_future\_\_ import annotations} had to be simplified,
\item SymPy's generic \texttt{simplify} was too expensive inside the recursion,
so I replaced most heavy simplifications by lighter \texttt{cancel} calls,
\item once those simplification bottlenecks were removed, the exact symbolic
five-point calculation completed.
\end{enumerate}

The chamber-fixed symbolic outputs were the turning point. For three different
sample chambers, I obtained:

\textbf{Chamber sample $(4,3,2)$.}
\[
\frac{A_5}{\ii}
=
-32\,a\,b^2 c^2\,
\frac{ab+ac+b^2+bc+c^2}{a+b+c}.
\]

\textbf{Chamber sample $(2,5/2,3)$.}
\[
\frac{A_5}{\ii}
=
-16\,a^5\,
\frac{ab+ac+b^2+bc+c^2}{a+b+c}.
\]

\textbf{Chamber sample $(5/2,2,3)$.}
\[
\frac{A_5}{\ii}
=
-16\,a\,b^2(2a^2-b^2)\,
\frac{ab+ac+b^2+bc+c^2}{a+b+c}.
\]

These are clearly not the same polynomial, so a naive chamber-independent
polynomial description was impossible. But the differences were structured, not
random. Each chamber changed only which quadratic pieces survived.

That suggested a new viewpoint:

\medskip
\textbf{Hypothesis E.} The answer is not best thought of as a rational function
first. It is better thought of as an inclusion--exclusion expression built from
truncated powers in the squared frequencies, whose polynomial reduction changes
from chamber to chamber.
\medskip

\subsection*{7. Recognition of the $n=5$ truncated-power pattern}

The five-point chamber outputs eventually collapsed to a single compact model.
Writing
\[
x=\omega_2^2,\qquad y=\omega_3^2,\qquad z=\omega_4^2,
\]
the chamber dependence is encoded by
\[
x^2-\pp{x-y}^2-\pp{x-z}^2+\pp{x-y-z}^2,
\qquad \pp{u}:=\max(u,0).
\]

So the five-point formula became
\[
A_5
=
16\,\ii\,\omega_1\omega_2
\left(
\omega_2^4
-\pp{\omega_2^2-\omega_3^2}^2
-\pp{\omega_2^2-\omega_4^2}^2
+\pp{\omega_2^2-\omega_3^2-\omega_4^2}^2
\right).
\]

This one formula immediately explained the sample data.

\textbf{Example 1: \texttt{freeW} = $\{1,2,3\}$.}
Here
\[
\omega = \{-4,1,2,3,-2\},
\qquad
\omega_2^2=1,\ \omega_3^2=4,\ \omega_4^2=9.
\]
All positive-part differences vanish except the first term, so
\[
A_5
=
16\,\ii\,(-4)(1)(1)
=
-64\,\ii.
\]

\textbf{Example 2: \texttt{freeW} = $\{2,1,3\}$.}
Here
\[
\omega = \{-7/2,2,1,3,-5/2\},
\qquad
\omega_2^2=4,\ \omega_3^2=1,\ \omega_4^2=9.
\]
Now
\[
\pp{4-1}=3,\qquad \pp{4-9}=0,\qquad \pp{4-1-9}=0,
\]
so
\[
A_5
=
16\,\ii\,\left(-\frac72\right)(2)\left(4^2-3^2\right)
=
-784\,\ii.
\]

\textbf{Example 3: \texttt{freeW} = $\{5,4,1\}$.}
Here
\[
\omega = \left\{-\frac{23}{5},5,4,1,-\frac{27}{5}\right\},
\]
so
\[
\omega_2^2=25,\qquad \omega_3^2=16,\qquad \omega_4^2=1.
\]
Then
\[
\pp{25-16}=9,\qquad \pp{25-1}=24,\qquad \pp{25-16-1}=8,
\]
and
\[
A_5
=
16\,\ii\,\left(-\frac{23}{5}\right)(5)\left(25^2-9^2-24^2+8^2\right)
=
-11776\,\ii.
\]

Once these checks worked, the five-point pattern looked stable.

\subsection*{8. Generalization to all $n$}

The five-point structure has an obvious higher-point completion. For
$n \ge 4$, define
\[
\pp{u}:=\max(u,0),
\]
and let the positive-sector interior legs be $3,\ldots,n-1$. The natural
general conjecture is
\[
A_n
=
\ii\,2^{n-1}\,\omega_1\omega_2
\sum_{S\subseteq\{3,\ldots,n-1\}}
(-1)^{|S|}
\pp{\omega_2^2-\sum_{j\in S}\omega_j^2}^{\,n-3}.
\]

Equivalently, if
\[
B_m(x;x_1,\ldots,x_m)
=
\sum_{S\subseteq\{1,\ldots,m\}}
(-1)^{|S|}
\pp{x-\sum_{i\in S}x_i}^{\,m},
\]
then
\[
A_n
=
\ii\,2^{n-1}\,\omega_1\omega_2\,
B_{n-3}\!\left(\omega_2^2;\omega_3^2,\ldots,\omega_{n-1}^2\right).
\]

This was the decisive general hypothesis:

\medskip
\textbf{Hypothesis F.} The entire two-minus sector is governed by a single
subset inclusion--exclusion formula in the squared positive-sector frequencies,
with chamber dependence carried only by the positive-part operation.
\medskip

This was a stronger and much more useful statement than any one-dimensional
rational fit. The only remaining question was whether it matched the original
Berends--Giele recursion beyond five points.

\subsection*{9. Exact verification at $n=5,6,7$}

To test Hypothesis F, I wrote \texttt{verify\_formula.wl}. It evaluates both
the supplied exact \texttt{BGAmplitude} and the conjectured closed form on
exact rational kinematics. The comparison points were chosen to mix easy
integers, half-integers, and less symmetric orderings.

The exact checks were:

\begin{center}
\small
\begin{longtable}{>{\raggedright\arraybackslash}p{0.07\textwidth} >{\raggedright\arraybackslash}p{0.23\textwidth} >{\raggedright\arraybackslash}p{0.28\textwidth} >{\raggedright\arraybackslash}p{0.17\textwidth} >{\raggedright\arraybackslash}p{0.17\textwidth}}
\toprule
$n$ & \texttt{freeW} & $\omega$ from \texttt{MakeKinematics} & $A_n$ from BG & closed form \\
\midrule
\endfirsthead
\toprule
$n$ & \texttt{freeW} & $\omega$ from \texttt{MakeKinematics} & $A_n$ from BG & closed form \\
\midrule
\endhead
5 & $\{1,2,3\}$ & $\{-4,1,2,3,-2\}$ & $-64\,\ii$ & $-64\,\ii$ \\
5 & $\{2,1,3\}$ & $\{-7/2,2,1,3,-5/2\}$ & $-784\,\ii$ & $-784\,\ii$ \\
5 & $\{5,4,1\}$ & $\{-23/5,5,4,1,-27/5\}$ & $-11776\,\ii$ & $-11776\,\ii$ \\
5 & $\{3/2,5/2,7/2\}$ & $\{-29/6,3/2,5/2,7/2,-8/3\}$ & $-(2349/4)\,\ii$ & $-(2349/4)\,\ii$ \\
5 & $\{5/2,7/2,3/2\}$ & $\{-43/10,5/2,7/2,3/2,-16/5\}$ & $-(15867/4)\,\ii$ & $-(15867/4)\,\ii$ \\
6 & $\{1,3/2,2,5/2\}$ & $\{-121/28,1,3/2,2,5/2,-75/28\}$ & $-(968/7)\,\ii$ & $-(968/7)\,\ii$ \\
6 & $\{5/2,2,3,7/2\}$ & $\{-70/11,5/2,2,3,7/2,-51/11\}$ & $-(1303400/11)\,\ii$ & $-(1303400/11)\,\ii$ \\
6 & $\{4,3,2,1\}$ & $\{-49/10,4,3,2,1,-51/10\}$ & $-(677376/5)\,\ii$ & $-(677376/5)\,\ii$ \\
6 & $\{7/2,5/2,3/2,1\}$ & $\{-139/34,7/2,5/2,3/2,1,-75/17\}$ & $-(656775/17)\,\ii$ & $-(656775/17)\,\ii$ \\
7 & $\{1,3/2,2,5/2,3\}$ & $\{-241/40,1,3/2,2,5/2,3,-159/40\}$ & $-(1928/5)\,\ii$ & $-(1928/5)\,\ii$ \\
7 & $\{4,3,5/2,2,1\}$ & $\{-321/50,4,3,5/2,2,1,-152/25\}$ & $-(426108561/50)\,\ii$ & $-(426108561/50)\,\ii$ \\
7 & $\{7/2,5/2,2,3/2,1\}$ & $\{-223/42,7/2,5/2,2,3/2,1,-109/21\}$ & $-1602701\,\ii$ & $-1602701\,\ii$ \\
7 & $\{5/2,7/2,3/2,9/4,3\}$ & $\{-29/4,5/2,7/2,3/2,9/4,3,-11/2\}$ & $-(12048407135/8192)\,\ii$ & $-(12048407135/8192)\,\ii$ \\
\bottomrule
\end{longtable}
\end{center}

Every one of those tests gave exact difference zero. In other words,
\[
A_n^{\text{BG}} - A_n^{\text{closed form}} = 0
\]
as an exact rational identity at each sampled point, not merely to floating
point tolerance.

There was one seven-point sample that I initially chose,
\[
\texttt{freeW}=\{5/2,7/2,3/2,2,3\},
\]
for which the raw BG recursion again returned \texttt{Indeterminate}. I did not
take that as evidence against the formula, because the same pathology had
already appeared at special lower-point singular kinematics. I replaced it by a
nearby generic rational sample,
\[
\texttt{freeW}=\{5/2,7/2,3/2,9/4,3\},
\]
and the exact agreement reappeared immediately. So the real lesson of those
exceptions is that the supplied recursive representation itself becomes
ill-defined on special exact channel loci, not that the closed formula fails.

\subsection*{10. Returning to four points}

Once the higher-point pattern was locked in, it was straightforward to return
to $n=4$. The same general expression gives
\[
A_4
=
8\,\ii\,\omega_1\omega_2\left(\pp{\omega_2^2}-\pp{\omega_2^2-\omega_3^2}\right).
\]
Since all squared frequencies are nonnegative on the real kinematic chart, this
simplifies to
\[
A_4
=
8\,\ii\,\omega_1\omega_2\,\min(\omega_2^2,\omega_3^2).
\]

That finite continuation gives, for example,
\[
\texttt{freeW}=\{2,3\}
\quad\Longrightarrow\quad
\omega=\{-3,2,3,-2\},
\quad
A_4=-192\,\ii,
\]
and
\[
\texttt{freeW}=\{5/2,7/2\}
\quad\Longrightarrow\quad
\omega=\{-7/2,5/2,7/2,-5/2\},
\quad
A_4=-\frac{875}{2}\,\ii.
\]

So the final four-point interpretation is:

\begin{enumerate}
\item the \emph{raw} exact BG recursion in \texttt{OnShellBG.m} returns
\texttt{Indeterminate} at four points in this sector because exact internal
zero-momentum channels are hit,
\item the \emph{closed formula inferred from higher points} has a finite and
simple four-point continuation.
\end{enumerate}

This resolved the earlier tension that produced Hypothesis B. The four-point
singularity belonged to the chosen recursive representation, not to the final
closed expression itself.

\subsection*{11. Final assessment}

Chronologically, the decisive changes of mind were:

\begin{enumerate}
\item direct symbolic four-point work looked too singular and too messy to be a
good entry point,
\item literal rational fitting on a one-parameter five-point slice did not
stabilize into a convincing global formula,
\item chamber-fixed exact symbolic formulas at five points exposed a much simpler
pattern than the prompt's original rational-language hint suggested,
\item that pattern was the truncated-power inclusion--exclusion combination in
squared frequencies,
\item the resulting all-$n$ conjecture matched exact BG data at every generic
five-, six-, and seven-point rational test I ran.
\end{enumerate}

So the final closed formula I ended up trusting was not reached by a single
global numerator/denominator fit. It was reached by:

\begin{enumerate}
\item collecting exact low-point data,
\item recognizing the right chamber language,
\item extracting local symbolic formulas,
\item identifying the inclusion--exclusion structure,
\item verifying the all-$n$ conjecture directly against the source recursion.
\end{enumerate}

The resulting artifacts in the output directory are:

\begin{enumerate}
\item \texttt{bg\_core.wl}: recursion core copied from the source file,
\item \texttt{symbolic\_bg.py}: chamber-fixed exact symbolic extractor,
\item \texttt{verify\_formula.wl}: exact verification script,
\item \texttt{verification.txt}: exact BG versus closed-form checks,
\item \texttt{result.md}: concise statement of the final formula.
\end{enumerate}

\end{document}
